
\section{\BookRecordIndex{T0233}临界点}\label{sec:2-4}
\zhindex{linjiedian}{临界点}

本节建立 \BookRecordIndex{T0003}$J$-全纯曲线的一些有用的局部性质。本节始终假设 $J$ 属于 $C^1$；于是由
注~\ref{rem:2-3-1}，$u$ 也属于 $C^1$。\BookRecordIndex{T0003}$J$-全纯曲线 $u:\Sigma\to M$ 的一个
\emph{\BookRecordIndex{T0233}临界点}是满足 $du(z)=0$ 的点 $z\in\Sigma$。\BookRecordIndex{T0233}临界点在 $u$ 下的像称为
\emph{\BookRecordIndex{T0234}临界值}。如果把像 $C:=u(\Sigma)\subset M$ 看成一条不带参数的 \BookRecordIndex{T0003}$J$-全纯曲线，
那么 $C$ 上的\BookRecordIndex{T0233}临界点就是 $u$ 的\BookRecordIndex{T0234}临界值。$C$ 上不是\BookRecordIndex{T0233}临界点的点称为\BookRecordIndex{T0405}正则点或\BookRecordIndex{T0121}非奇异点。
在可积情形，众所周知，非常值全纯曲线的\BookRecordIndex{T0233}临界点是孤立的。下一个引理断言，对任意\BookRecordIndex{T0088}殆复
结构也有同样结论。我们给出两个证明。

\begin{lemma}\label{lem:2-4-1}
设 $\Sigma$ 是紧致无边界黎曼曲面，$J$ 是光滑流形 $M$ 上的 $C^1$ \BookRecordIndex{T0088}殆复结构，而
$u:\Sigma\to M$ 是一条非常值的 $C^1$ 类 \BookRecordIndex{T0003}$J$-全纯曲线。则\BookRecordIndex{T0234}临界值的原像集
\[
  X:=u^{-1}\bigl(\{u(z)\mid z\in\Sigma,\ du(z)=0\}\bigr)
\]
是有限集。此外，对每个 $x\in M$，$u^{-1}(x)$ 都是有限集。
\end{lemma}

\begin{proof}[证明 1]
在第一个证明中，假设 $J$ 光滑。只需证明\BookRecordIndex{T0233}临界点是孤立的，因此可以局部工作。于是设
$\Omega\subset\mathbb C$ 是零点的一个开邻域，
$u:\Omega\to\mathbb C^n$ 关于某个\BookRecordIndex{T0088}殆复结构
$J:\mathbb C^n\to\GL(2n,\mathbb R)$ 是 $J$-全纯的，并且
\[
  u(0)=0,\qquad du(0)=0,\qquad u\not\equiv0,\qquad J(0)=J_0.
\]
写 $z=s+it$。因为 $u$ 不是常值，由推论~\ref{cor:2-3-3}，$u(z)$ 在 $z=0$ 处的
无穷阶 jet 必不为零。因此存在整数 $\ell\ge2$，使得
\[
  u(z)=O(|z|^\ell),\qquad u(z)\ne O(|z|^{\ell+1}).
\]
这蕴含
\[
  J(u(z))=J_0+O(|z|^\ell).
\]
现在把方程
\[
  \partial_su+J(u)\partial_tu=0
\]
的左边作到 $\ell-1$ 阶的 Taylor 展开，得到
\[
  \partial_sT_\ell(u)+J_0\partial_tT_\ell(u)=0.
\]
这里 $T_\ell(u):\mathbb C\to\mathbb C^n$ 表示 $u$ 到 $\ell$ 阶为止的 Taylor 展开。
由此 $T_\ell(u)$ 是全纯函数，并且存在非零向量 $a\in\mathbb C^n$，使得
\[
  u(z)=az^\ell+O(|z|^{\ell+1}),\qquad
  \partial_su(z)=\ell az^{\ell-1}+O(|z|^\ell).
\]
因而，当 $\varepsilon>0$ 充分小时，
\[
  0<|z|\le\varepsilon
  \qquad\Longrightarrow\qquad
  u(z)\ne0,\quad du(z)\ne0.
\]
所以\BookRecordIndex{T0233}临界点是孤立的；又因 $\Sigma$ 紧，$u$ 的\BookRecordIndex{T0233}临界点集是有限集。还可\BookRecordIndex{T0326}推出，对每个
$p\in M$，$u^{-1}(p)$ 都是有限集。因此 $X$ 是有限集。
\end{proof}

\begin{proof}[证明 2]
设 $J:\mathbb C^n\to\GL(2n,\mathbb R)$ 是 $C^1$ 类\BookRecordIndex{T0088}殆复结构，
$\Omega\subset\mathbb C$ 是开集，并设 $u$ 是这样一条 \BookRecordIndex{T0003}$J$-全纯曲线。则 $u$ 满足
\eqref{eq:2-3-4}，其中 $J(z)=J(u(z))$ 属于 $C^1$ 且 $C(z)=0$。因此由
定理~\ref{thm:2-3-5}，对每个 $p\in M$，$u^{-1}(p)$ 都是离散集；回到紧曲面
$\Sigma$，可知这些原像集是有限集。\sourcefootnote{原书在局部开集 $\Omega$ 的表述中
直接称 $u^{-1}(p)$ 为有限集；定理 2.3.5 在这里给出的是局部孤立性，而有限性还须使用
全局定义域 $\Sigma$ 的紧致性。}此外，对
$C=0$ 的方程
\eqref{eq:2-3-4} 求导可知，函数
$v:=\partial_su\in W^{1,p}(\Omega,\mathbb C^n)$ 满足
\[
  \partial_sv(z)+J(z)\partial_tv(z)
  +\bigl(\partial_sJ(z)\bigr)J(z)v(z)=0.
\]
把定理~\ref{thm:2-3-5} 应用于 $v$，可知它的零点都是孤立的。回到 $\Sigma$，因此
$u$ 的\BookRecordIndex{T0233}临界点集是有限集，$X$ 也同样是有限集。
\end{proof}

下面说明怎样在 \BookRecordIndex{T0003}$J$-全纯曲线的\BookRecordIndex{T0405}正则点附近选取良好的坐标。

\begin{lemma}\label{lem:2-4-2}
设 $\ell\ge2$，$J$ 是光滑流形 $M$ 上的 $C^\ell$ \BookRecordIndex{T0088}殆复结构。设
$\Omega\subset\mathbb C$ 是零点的一个开邻域，且
$u:\Omega\to M$ 是满足 $du(0)\ne0$ 的局部 \BookRecordIndex{T0003}$J$-全纯曲线。则存在定义在 $u(0)$
的一个开邻域上的 $C^{\ell-1}$ 坐标卡 $\psi:U\to\mathbb C^n$，使得
\[
  \psi\circ u(z)=(z,0,\ldots,0),\qquad
  d\psi(u(z))J(u(z))=J_0d\psi(u(z))
\]
对所有 $z\in\Omega\cap u^{-1}(U)$ 成立。
\end{lemma}

\begin{proof}
由定理 B.4.1 和注 B.4.3，对每个 $p<\infty$，$u$ 都属于 $W^{\ell+1,p}$；因此由
\BookRecordIndex{T0318}索伯列夫嵌入定理 B.1.11，$u$ 属于 $C^\ell$。写
$z=s+it\in\Omega$，并写
$w=(w_1,w_2,\ldots,w_n)\in\mathbb C^n$，其中 $w_j=x_j+iy_j$。必要时缩小
$\Omega$，并选取丛 $u^*TM$ 的一个复 $C^{\ell-1}$ \BookRecordIndex{T0067}标架，使得
\[
  Z_1(z),\ldots,Z_n(z)\in T_{u(z)}M,\qquad
  Z_1=\frac{\partial u}{\partial s}.
\]
定义 $\phi:\Omega\times\mathbb C^{n-1}\to M$ 为
\[
  \phi(w_1,\ldots,w_n)
  =\exp_{u(w_1)}\left(
    \sum_{j=2}^n x_jZ_j(w_1)
    +\sum_{j=2}^n y_jJ(u(w_1))Z_j(w_1)
  \right).
\]
于是 $\phi$ 把 $\mathbb C^n$ 中零点的一个邻域 $V$ 以 $C^{\ell-1}$ 微分同胚映到
$M$ 中 $u(0)$ 的一个邻域 $U$。它满足
$\phi(w_1,0,\ldots,0)=u(w_1)$，并且在所有
$w=(w_1,0,\ldots,0)$ 处都有
\[
  \frac{\partial\phi}{\partial x_j}
  +J(\phi)\frac{\partial\phi}{\partial y_j}=0,
  \qquad j=1,\ldots,n.
\]
因此逆映射 $\psi:=\phi^{-1}:U\to V$ 具有所要求的性质。
\end{proof}

下面开始研究两条不同的 \BookRecordIndex{T0003}$J$-全纯曲线之间的交点。与此有关的最重要结果出现在四维，并且例如
在 McDuff \srcite{256, 264} 中有所讨论。目前先证明一个在所有维数都成立的有用结果：
两条不同的
\BookRecordIndex{T0003}$J$-全纯曲线 $u:\Sigma\to M$ 和 $u':\Sigma'\to M$ 的交点，只可能在同时为两条曲线
$C=u(\Sigma)$ 与 $C'=u'(\Sigma')$ 的\BookRecordIndex{T0233}临界点处聚集。\footnote{事实上，任意两条不同的闭
\BookRecordIndex{T0003}$J$-全纯曲线都只有有限多个交点（见 Sikorav \srcite{379} 以及 Lazzarini 所写的附录 E）。
不过，这一点证明起来更困难，而且在我们的多数应用中并不需要。这里的结果取自 McDuff
\srcite{254}。} 对局部 \BookRecordIndex{T0003}$J$-全纯曲线，这个陈述可以改写如下。

\begin{lemma}\label{lem:2-4-3}
设 $J$ 是光滑流形 $M$ 上的 $C^2$ \BookRecordIndex{T0088}殆复结构，$\Omega\subset\mathbb C$ 是零点的一个开
邻域，而 $u,v:\Omega\to M$ 是满足
\[
  u(0)=v(0),\qquad du(0)\ne0
\]
的 \BookRecordIndex{T0003}$J$-全纯曲线。再假设存在点列 $z_\nu,\zeta_\nu\in\Omega$，使得
\[
  u(z_\nu)=v(\zeta_\nu),\qquad
  \lim_{\nu\to\infty}z_\nu
  =\lim_{\nu\to\infty}\zeta_\nu=0,
  \qquad \zeta_\nu\ne0.
\]
则在零点的某个邻域 $B_\varepsilon(0)$ 上存在全纯函数
$\phi:B_\varepsilon(0)\to\Omega$，使得 $\phi(0)=0$ 且
\[
  v=u\circ\phi.
\]
\end{lemma}

\begin{proof}[证明 1]
在这个证明中，假设 $J$ 光滑，并使用 Aronszajn 定理。由引理~\ref{lem:2-4-2}，不失
一般性，可以假设 $M=\mathbb C^n$，
$J:\mathbb C^n\to\GL(2n,\mathbb R)$ 是光滑\BookRecordIndex{T0088}殆复结构，并且
\[
  u(z)=(z,0),\qquad J(w_1,0)=i,
\]
其中 $w=(w_1,\widetilde w)$ 且 $\widetilde w\in\mathbb C^{n-1}$。写
$v(z)=:(v_1(z),\widetilde v(z))$。

先证明 $\widetilde v$ 在 $z=0$ 处的无穷阶 jet 必须消失。否则，存在整数
$\ell\ge0$，使得
$\widetilde v(z)=O(|z|^\ell)$ 且
$\widetilde v(z)\ne O(|z|^{\ell+1})$。引理的假设蕴含 $\ell\ge1$，从而
$J(v(z))=J_0+O(|z|^\ell)$。如引理~\ref{lem:2-4-1} 的证明那样，把方程
$\partial_sv+J\partial_tv=0$ 左边作到 $\ell-1$ 阶的 Taylor 展开，得到
$T_\ell(v)$ 全纯。因此
\[
  v_1(z)=p(z)+O(|z|^{\ell+1}),\qquad
  \widetilde v(z)=\widetilde a z^\ell+O(|z|^{\ell+1}),
\]
其中 $p(z)$ 是一个 $\ell$ 阶多项式，
$0\ne\widetilde a\in\mathbb C^{n-1}$。这说明在零点的某个邻域内，每个非零点
$\zeta$ 都满足 $\widetilde v(\zeta)\ne0$，从而
$v(\zeta)\notin\operatorname{im}u$，与引理的假设矛盾。因此已经证明
$\widetilde v$ 在 $z=0$ 处的无穷阶 jet 消失。

再证明 $\widetilde v(z)\equiv0$。注意，因为在轴 $\widetilde w=0$ 上有 $J=J_0$，
所以对所有 $w_1$，
\[
  \frac{\partial J(w_1,0)}{\partial x_1}
  =\frac{\partial J(w_1,0)}{\partial y_1}=0.
\]
因而
\[
  \left|\frac{\partial J(w)}{\partial x_1}\right|
  +\left|\frac{\partial J(w)}{\partial y_1}\right|
  \le c|\widetilde w|.
\]
利用 \eqref{eq:2-3-2}，得到不等式
\[
  |\Delta\widetilde v|
  \le c\bigl(|\widetilde v|+|\partial_s\widetilde v|
  +|\partial_t\widetilde v|\bigr).
\]
于是由 Aronszajn 定理，$\widetilde v\equiv0$。因此
$v(z)=(v_1(z),0)$，这蕴含 $\phi:=v_1$ 是全纯的，并且对 $z\in B_\delta$ 有
$v(z)=u(\phi(z))$，正是所需结论。
\end{proof}

\begin{proof}[证明 2]
在这个证明中，考虑一般情形：$J$ 属于 $C^\ell$，其中 $\ell\ge2$ 是某个常数，并使用
\BookRecordIndex{T0013}Carleman 相似原理。由引理~\ref{lem:2-4-2}，可以假设
$J:\mathbb C^n\to\GL(2n,\mathbb R)$、$u(z)=(z,0)$ 以及
$v=(v_1,\widetilde v)$ 如第一个证明中所述，只是现在 $J$ 和 $v$ 仅连续可微。我们断言，
存在连续函数
$\widetilde C:\Omega\to\End_{\mathbb R}(\mathbb C^{n-1})$，使得
\[
  \partial_s\widetilde v+i\partial_t\widetilde v
  +\widetilde C\widetilde v=0.
\]
这个方程可以写成
\[
  \widetilde C\widetilde v
  =\widetilde\pi\bigl((J(v)-i)\partial_tv\bigr),
\]
其中 $\widetilde\pi:\mathbb C^n\to\mathbb C^{n-1}$ 表示投影
$\widetilde\pi(w_1,\widetilde w):=\widetilde w$。因为 $J(v_1,0)=i$，函数
$\widetilde C$ 由下式给出：
\[
  \widetilde C(z)\widetilde\xi
  :=\widetilde\pi\left(
    \left(
      \int_0^1
      \left.\frac{d}{ds}\right|_{s=0}
      J\bigl(v_1(z),\tau\widetilde v(z)+s\widetilde\xi\bigr)\,d\tau
    \right)\partial_tv(z)
  \right)
\]
对 $z\in\Omega$ 和 $\widetilde\xi\in\mathbb C^{n-1}$ 成立。注意，
$\widetilde C$ 是连续的。

现在由定理~\ref{thm:2-3-5}，存在常数 $\delta>0$、函数
\[
  \widetilde\Phi\in W^{1,p}\bigl(B_\delta,\GL(n-1,\mathbb C)\bigr)
\]
以及全纯函数 $\widetilde\sigma:B_\delta\to\mathbb C^{n-1}$，使得
\[
  \widetilde v(z)=\widetilde\Phi(z)\widetilde\sigma(z),
  \qquad \widetilde\sigma(0)=0.
\]
由假设，存在点列 $0\ne\zeta_\nu\to0$，满足
$\widetilde v(\zeta_\nu)=0$。因此，当 $\nu$ 充分大时，
$\widetilde\sigma(\zeta_\nu)=0$，这蕴含 $\widetilde v\equiv0$。
\end{proof}

下面证明前一个引理的全局版本；这是点集拓扑中的一个简单习题。第 3.4 节讨论\BookRecordIndex{T0283}求值映射
的\BookRecordIndex{T0166}横截性时将会用到它。

\begin{proposition}\label{prop:2-4-4}
设 $J$ 是 $M$ 上的 $C^2$ \BookRecordIndex{T0088}殆复结构，$\Sigma_0,\Sigma_1$ 是紧致、连通且无边界的黎曼曲面。
设 $u_0:\Sigma_0\to M$ 和 $u_1:\Sigma_1\to M$ 是 \BookRecordIndex{T0003}$J$-全纯曲线，并且
$u_0(\Sigma_0)\ne u_1(\Sigma_1)$，$u_0$ 不是常值。则集合
$u_0^{-1}(u_1(\Sigma_1))$ 至多可数，并且只能在 $u_0$ 的\BookRecordIndex{T0233}临界点处聚集。
\end{proposition}

\begin{proof}
如果 $u_1$ 是常值，则由引理~\ref{lem:2-4-1}，
$u_0^{-1}(u_1(\Sigma_1))$ 是有限集，因而无须证明。以下假设 $u_1$ 不是常值。对
$i=0,1$，以 $X_i\subset\Sigma_i$ 表示 $u_i$ 的\BookRecordIndex{T0233}临界点集。由
引理~\ref{lem:2-4-1}，这些集合都是有限集。我们证明，对每个
$z_0\in\Sigma_0\setminus X_0$，下列条件等价：
\begin{enumerate}[label=(\roman*),leftmargin=3em]
  \item 存在 $z_0$ 的一个邻域 $U_0\subset\Sigma_0$，使得
    $u_0(U_0)\subset u_1(\Sigma_1)$；
  \item 存在点列
    $z_\nu\in u_0^{-1}(u_1(\Sigma_1))\setminus\{z_0\}$ 收敛到 $z_0$。
\end{enumerate}
条件 (i) 显然蕴含 (ii)。反过来，假设 (ii)，并选取点列 $\zeta_\nu\in\Sigma_1$，
使得对每个 $\nu$ 都有 $u_1(\zeta_\nu)=u_0(z_\nu)$。必要时取子列，可以假设
$\zeta_\nu$ 收敛到 $\zeta_0$，于是
$u_1(\zeta_0)=u_0(z_0)$。由引理~\ref{lem:2-4-2}，$u_0$ 在 $z_0$ 附近
单射\sourcefootnote{原书此处仅由 $z_\nu\ne z_0$ 直接\BookRecordIndex{T0326}推出
$u_0(z_\nu)\ne u_0(z_0)$；这一推论对一般映射并不成立。这里利用
$du_0(z_0)\ne0$ 以及引理 2.4.2 给出的局部单射性补足此步。}；结合
$z_\nu\ne z_0$ 和 $z_\nu\to z_0$ 可知，对充分大的 $\nu$，
$u_0(z_\nu)\ne u_0(z_0)$，因而 $\zeta_\nu\ne\zeta_0$。因为 $z_0$
不是 $u_0$ 的\BookRecordIndex{T0233}临界点，由引理~\ref{lem:2-4-3}，存在 $\zeta_0$ 的一个邻域
$U_1\subset\Sigma_1$ 以及全纯映射 $\phi:U_1\to\Sigma_0$，使得
$\phi(\zeta_0)=z_0$ 且 $u_1|_{U_1}=u_0\circ\phi$。因为 $u_1$ 不是常值，$\phi$ 也不是
常值。因此 $U_0:=\phi(U_1)$ 是 $z_0$ 的一个邻域，并且
$u_0(U_0)\subset u_1(\Sigma_1)$。这就证明了 (i)。

现在令 $W_0\subset\Sigma_0\setminus X_0$ 为所有满足 (i)、因而也满足 (ii) 的点
$z_0$ 的集合。由 (i)，这个集合是开的；由 (ii)，它在
$\Sigma_0\setminus X_0$ 中相对闭。我们必须证明 $W_0=\varnothing$。反设
$W_0\ne\varnothing$。于是 $W_0=\Sigma_0\setminus X_0$，因而
$u_0(\Sigma_0)\subset u_1(\Sigma_1)$。这又蕴含，以互换 $\Sigma_0$ 与 $\Sigma_1$
角色的同样方式定义的集合 $W_1\subset\Sigma_1\setminus X_1$ 也非空。（请证明！）
因此，同样的论证表明 $u_1(\Sigma_1)\subset u_0(\Sigma_0)$，从而
$u_0(\Sigma_0)=u_1(\Sigma_1)$。这与假设矛盾，所以正如所断言的，
$W_0=\varnothing$。命题的结论现在由 $W_0$ 的定义得出。
\end{proof}
